\section{Training data and configuration}
\label{app:training}

\paragraph{Data construction and splitting.}
The generators construct game states and legal histories, replay observations,
and solve the resulting decisions without using an LLM to supply labels.
Initial commitments specify the starting state; voluntary partner actions
provide behavioral evidence under the reference policy. Investigation
branches expose truthful preference feedback and the subsequent legal
decisions. The paired bank renders each canonical decision as direct action
selection, inference, and conditional planning. Its split groups source
families, including all views of a decision and related histories. This is a
split within the constructed bank: some families from earlier development
validation sets were reassigned to training, so it does not establish
separation from every historical evaluation suite.

\paragraph{Planning labels and dataset versions.}
For the paired bank, planning prompts supply the current state and a
qualitative oracle judgment while omitting behavioral history. The label
audit distinguishes actions that are acceptable, unacceptable, or ambiguous
given that judgment; ambiguous labels receive no task supervision. Of the
373 training cases, 314 support planning training, 29 have no negative
planning label and are reserved for evaluation, and 30 are quarantined.
The corresponding validation counts are 106, 4, and 14. The historical
\textsc{GameSlice-IP} checkpoint uses an earlier dataset with 294 inference
and 300 planning items. These earlier planning prompts use a mixture of
interfaces, as do the planning items retained for compact joint training;
they should not be identified with the uniform intervention interface in
Section~\ref{sec:diagnose}.

\paragraph{Selecting the compact joint dataset.}
The 24-item joint set contains eight examples per objective, selected from
training data using archived learning histories. For direct action selection
and inference, the main selection rule requires at least four prior
exposures, two groups with nonzero task advantages, mean accuracy of at least
$50\%$ over the final two exposures, and an improvement of at least
$12.5$ percentage points over the first two. One inference item has only
three prior exposures and is retained as a documented exception. The eight
planning items come from an earlier compact selection. Original prompts
and teachers are preserved. Neither CalBench results nor validation
performance enters item selection. The resulting set is frozen before
training: two deterministic halves alternate across 40 updates, giving
each item 20 groups of eight responses. Because selection uses prior
learning performance, this run measures training on a targeted subset,
rather than an isolated change in objective or exposure count.

\begingroup
\small
\renewcommand{\arraystretch}{1.12}
\setlength{\tabcolsep}{5pt}
\begin{longtable}{@{}p{0.34\textwidth}p{0.62\textwidth}@{}}
\caption{Data and training settings. Shared optimizer and decoding settings
are verified in the archived configurations for the selected
\textsc{GameSlice-A}, \textsc{GameSlice-IP}, and \textsc{GameSlice-Joint}
checkpoints. Recipe-specific rows distinguish nominal budgets from consumed
response tokens. A dash denotes a quantity not yet verified for the selected
checkpoint.}
\label{tab:training-configuration}\\
\hline
Parameter & Setting \\
\hline
\endfirsthead
\multicolumn{2}{l}{\tablename~\thetable\ (continued)}\\
\hline
Parameter & Setting \\
\hline
\endhead
\hline
\multicolumn{2}{r}{Continued on next page}\\
\endfoot
\hline
\endlastfoot
\multicolumn{2}{l}{\textbf{Shared slice-training configuration}}\\
Initial model / KL reference & Qwen3-4B-Instruct-2507 / frozen initial model \\
Precision & BF16; gradient accumulation and reduction in FP32 \\
Optimizer & Adam, Megatron backend \\
Initial learning rate & $10^{-6}$ \\
Adam coefficients; weight decay & $(0.9,0.95)$; $0$ \\
Warmup & $0$ steps \\
Policy updates per rollout batch & $1$ epoch \\
Policy-ratio clipping & $0.2$ \\
KL loss coefficient & $0.01$ \\
Loss aggregation & Mean over tokens within each response, then over responses \\
Sequence / prompt limit & $4096$ / $3072$ tokens \\
Maximum generated response & $1024$ tokens per call \\
Sampling & Temperature $1$; top-$p=1$; top-$k$ disabled \\
Repetition penalty & $1$ \\
Responses per question group & $8$ \\
Invalid or truncated output & Independent negative advantage, coefficient $0.2$ \\
Random seed & $42$ \\
Training hardware layout & One node, two GPUs; tensor parallelism $2$, pipeline parallelism $1$ \\
Microbatch size & $1$ per device \\
Memory settings & Sequence parallelism; full activation recomputation \\
Generation backend & vLLM; tensor parallelism $1$ per replica \\
Validation / checkpoint interval & Every $10$ updates \\
\hline
\multicolumn{2}{l}{\textbf{Datasets}}\\
Paired bank & $497$ canonical cases: $373$ training, $124$ validation; $175$ parent packages \\
Paired views & One direct-action, inference, and planning view per case; planning eligibility audited separately \\
Planning-eligible cases & $314$ training; $106$ validation \\
Historical inference--planning set & $594$ training items: $294$ inference, $300$ planning; $495$ validation items \\
Compact joint set & $24$ training items: $8$ direct-action, $8$ inference, $8$ planning \\
\hline
\multicolumn{2}{l}{\textbf{Selected checkpoints and realized training dose}}\\
\textsc{GameSlice-A} & O-99; $100$ updates on paired-bank direct-action items \\
Action-only budget / consumed & $6{,}553{,}600$ / $6{,}554{,}714$ response tokens \\
Action-only collection target & $65{,}536$ response tokens per update; complete groups at the boundary \\
\textsc{GameSlice-IP} & BP-99; $100$ updates on the historical inference--planning set \\
Inference--planning objective weights & $1/2$ inference, $1/2$ planning \\
Inference--planning budget / consumed & $6{,}553{,}600$ / $2{,}516{,}803$ response tokens at the selected checkpoint \\
Inference--planning learning rate & Configured constant rate of $10^{-6}$ \\
\textsc{GameSlice-Joint} & micro24-39; $40$ updates on the compact joint set \\
Joint batch composition & $4$ groups per objective; $12\times8=96$ responses per update \\
Joint objective weights & $1/3$ per objective \\
Joint exposure per item & $20$ groups; $160$ responses \\
Joint token cap / consumed & $3{,}932{,}160$ / $1{,}782{,}494$ response tokens \\
Action-only / joint learning rate & Token-based cosine decay from $10^{-6}$ toward $10^{-7}$, parameterized by each run's token budget \\
SP checkpoint & SP-41; complete-game self-play with own terminal utility \\
SP-41 cumulative training dose & -- \\
\end{longtable}
\endgroup

The response-token count measures generated tokens entering training,
including retries; it is distinct from the number of unique examples or
updates. The compact run terminates at its update limit before reaching its
token cap. Its count follows the completed checkpoint lineage and excludes
discarded work after an interrupted save. The selected SP-41 evaluation
artifact identifies the model export, but its cumulative training dose
remains to be linked to the training archive. Earlier SP configurations are
therefore not substituted for that checkpoint. The O and D labels in the
historical diagnosis tables denote action-only and joint training,
respectively; those earlier D checkpoints are distinct from micro24-39.
